\documentclass{article}

% if you need to pass options to natbib, use, e.g.:
%\PassOptionsToPackage{numbers, compress}{natbib}
% before loading neurips_2026

% The authors should use one of these tracks.
% Before accepting by the NeurIPS conference, select one of the options below.
% 0. "default" for submission
\usepackage[nonatbib]{neurips_2026}
\usepackage{neurips_2026}
% the "default" option is equal to the "main" option, which is used for the Main Track with double-blind reviewing.
% 1. "main" option is used for the Main Track
%  \usepackage[main]{neurips_2026}
% 2. "position" option is used for the Position Paper Track
%  \usepackage[position]{neurips_2026}
% 3. "eandd" option is used for the Evaluations & Datasets Track
 % \usepackage[eandd]{neurips_2026}
 % if you need to opt-in for a single-blind submission in the E&D track:
 %\usepackage[eandd, nonanonymous]{neurips_2026}
% 4. "creativeai" option is used for the Creative AI Track
%  \usepackage[creativeai]{neurips_2026}
% 5. "sglblindworkshop" option is used for the Workshop with single-blind reviewing
 % \usepackage[sglblindworkshop]{neurips_2026}
% 6. "dblblindworkshop" option is used for the Workshop with double-blind reviewing
%  \usepackage[dblblindworkshop]{neurips_2026}

% After being accepted, the authors should add "final" behind the track to compile a camera-ready version.
% 1. Main Track
 % \usepackage[main, final]{neurips_2026}
% 2. Position Paper Track
%  \usepackage[position, final]{neurips_2026}
% 3. Evaluations & Datasets Track
 % \usepackage[eandd, final]{neurips_2026}
% 4. Creative AI Track
%  \usepackage[creativeai, final]{neurips_2026}
% 5. Workshop with single-blind reviewing
%  \usepackage[sglblindworkshop, final]{neurips_2026}
% 6. Workshop with double-blind reviewing
%  \usepackage[dblblindworkshop, final]{neurips_2026}
% Note. For the workshop paper template, both \title{} and \workshoptitle{} are required, with the former indicating the paper title shown in the title and the latter indicating the workshop title displayed in the footnote.
% For workshops (5., 6.), the authors should add the name of the workshop, "\workshoptitle" command is used to set the workshop title.
% \workshoptitle{WORKSHOP TITLE}

% "preprint" option is used for arXiv or other preprint submissions
 % \usepackage[preprint]{neurips_2026}

% to avoid loading the natbib package, add option nonatbib:
\usepackage[nonatbib]{neurips_2026}
\usepackage{mathtools}
\usepackage{amssymb}
\usepackage{natbib} 
\usepackage{amsmath}
\usepackage[utf8]{inputenc} % allow utf-8 input
\usepackage[T1]{fontenc}    % use 8-bit T1 fonts
\usepackage{hyperref}       % hyperlinks
\usepackage{url}            % simple URL typesetting
\usepackage{booktabs}       % professional-quality tables
\usepackage{amsfonts}       % blackboard math symbols
\usepackage{nicefrac}       % compact symbols for 1/2, etc.
\usepackage{microtype}      % microtypography
\usepackage{xcolor}         % colors
\usepackage{algorithm} % For algorithm environment
\usepackage{algorithmic}
\usepackage{multirow}
\usepackage{amsthm}

% Define the assumption style
\newtheorem{assumption}{Assumption}
\newtheorem{theorem}{Theorem}
\newtheorem{lemma}[theorem]{Lemma}
\newtheorem{corollary}[theorem]{Corollary}
\newtheorem{proposition}[theorem]{Proposition}
% Note. For the workshop paper template, both \title{} and \workshoptitle{} are required, with the former indicating the paper title shown in the title and the latter indicating the workshop title displayed in the footnote. 
\title{Flow-based Framework Solving Multiple Processes PDEs}


% The \author macro works with any number of authors. There are two commands
% used to separate the names and addresses of multiple authors: \And and \AND.
%
% Using \And between authors leaves it to LaTeX to determine where to break the
% lines. Using \AND forces a line break at that point. So, if LaTeX puts 3 of 4
% authors names on the first line, and the last on the second line, try using
% \AND instead of \And before the third author name.


\author{%
   \\
  % examples of more authors
  % \And
  % Coauthor \\
  % Affiliation \\
  % Address \\
  % \texttt{email} \\
  % \AND
  % Coauthor \\
  % Affiliation \\
  % Address \\
  % \texttt{email} \\
  % \And
  % Coauthor \\
  % Affiliation \\
  % Address \\
  % \texttt{email} \\
  % \And
  % Coauthor \\
  % Affiliation \\
  % Address \\
  % \texttt{email} \\
}


\begin{document}


\maketitle


\begin{abstract}
  This paper proposes a novel method for solving multiple processes PDEs using a few-step flow matching with multiple vector fields, which correspond to multiple processes. Each vector field is modeled to depend not only on the corresponding process but also on the other processes, which reflect the interaction between different processes in a complex physics system.
\end{abstract}



\section{Introduction}
[TODO]
\section{Related Works}
[TODO]
\section{Method}
\label{sec:method}
\subsection{Problem Setting}
In this section, we provide a simple example of the coupled processes system as well as the multiple processes system. We begin with the coupled processes system with coupled kernels $\kappa_1$ and $\kappa_2$, coupled variable $v_1(x, t)$ and $v_2(x, t)$ given initial condition $v_1(x, 0)$ and $v_2(x, 0)$. Given $\mathcal{U}$ and $\mathcal{A}$ as the initial condition space and the prediction space, let $T$ denote the generic operator such that $T: \mathcal{U} \to \mathcal{A}$. Without the knowledge of how these two variables are coupled, to solve for $v_1(x, \tau)$ and $v_2(x, \tau)$, we need two operator $T_1$ and $T_2$ such that 
\begin{equation}
    v_1(x, \tau) = T_1 v_1(x, 0) = \int_{D} \kappa_1\big(x,y, v_1(x,0), v_2(x,0), v_1(y,0), v_2(y,0), \kappa_2\big) v_1(y,0) dy
\end{equation}
\begin{equation}
    v_2(x, \tau) = T_2 v_2(x, 0) = \int_{D} \kappa_2\big(x,y, v_1(x,0), v_2(x,0), v_1(y,0), v_2(y,0), \kappa_1\big) v_2(y,0) dy
\end{equation}
where $D \subset \mathbb{R}^{d}$ is a bounded domain. The interacting kernels cannot be directly solved without considering the interference from the other kernel. We can rewrite as 
\begin{equation}
    v_1(x, \tau) = T_1 [v_1(x, 0), v_2(x, 0); \kappa_1, \kappa_2] 
\end{equation}
\begin{equation}
    v_2(x, \tau) = T_2 [v_1(x, 0), v_2(x, 0); \kappa_1, \kappa_2] 
\end{equation}
Using the above formulations, we can easily extend to the multiple-process system, which includes $n$ different processes, such that, with $i \in [1, n]$, the process $i-$ th can be expressed as: 
\begin{equation}
    v_i(x, \tau) = T_i[v_1(x, 0), ...,v_n(x, 0); \kappa_1,..., \kappa_n]
\end{equation}
Our idea is to model each process operator $T_i$ by a flow model $\Psi_i$ with average velocity field from $r$ to $t$ as $u^{\theta}_i(z^t_1,..z^t_n; r,t )$, where we let $z^t_i (\cdot)$ denote the \emph{flow-matching state} at denoising time $t \in (0,1)$ of the physical process $i$-th such that the endpoints correspond to its initial and final physical time:  
% as immediate representation of flow $i$-th at time $t \in[0,1]$, such that 
\begin{equation}
    z^1_i(x) = v_i(x, 0)
\end{equation}
\begin{equation}
    z^0_i(x) = v_i(x, \tau). 
\end{equation}
Following standard flow matching, we define the regression target average velocity from flow matching time $0$ to $1$ of the $i-$th process as: 
\begin{equation}
    u_i^{\text{tar}}(z^1_1(x),..z^1_n(x); 0,1 ) = v_i(x, 0) -  v_i(x, \tau) \label{eq:global_tar_velocity}
\end{equation}
\subsection{Average Velocity with Coupled PDEs Processes}
\label{sec:average_velocity_couple}
In this section, we first derive the ground truth average velocity fields in the setting of two processes $v_1(x, t)$ and $v_2(x, t)$ as \citep{geng2025mean}. Let $u^{\text{ins}}_i(z^t_1(x), z^t_2(x), t )$ be instant velocity at time $t$ of process $i$-th, we have the average velocity $u_1$ from $r$ to $t$ as 
\begin{equation}
    (t- r) u_1(z^t_1(x), z^t_2(x), r, t) = \int_{r}^t u^{\text{ins}}_1(z^\tau_1(x), z^\tau_2(x), \tau ) d\tau
\end{equation}
Differentiating both sides with respect to $t$, we have:
\begin{equation}
    u_1(z^t_1(x), z^t_2(x), r, t) + (t -r) \frac{d}{dt}  u_1(z^t_1(x), z^t_2(x), r, t)= u^{\text{ins}}_1(z^t_1(x), z^t_2(x), t ) \label{eq:def_average_v}
\end{equation}
Note that $\frac{dr}{dt} = 0$ and $u^{\text{ins}}_1(z^t_1(x), z^t_2(x), t ) = \frac{d}{dt} (z^t_1(x))$ and $u^{\text{ins}}_2(z^t_1(x), z^t_2(x), t ) = \frac{d}{dt} (z^t_2(x))$, we have the time derivative of average velocity $u_1(z^t_1(x), z^t_2(x), r, t)$ as: 
\begin{align}
    \frac{d}{dt}  u_1(z^t_1(x), z^t_2(x), r, t) &=\partial_t u_1(z^t_1(x), z^t_2(x), r, t) \label{eq:time_der_coupled} \\ \nonumber
    &+ u^{\text{ins}}_1( z^t_1(x), z^t_2(x), t ) \partial_{z_1^t} u_1 (z^t_1(x), z^t_2(x), r, t) \\ \nonumber
    &+u^{\text{ins}}_2(z^t_1(x), z^t_2(x), t ) \partial_{z_2^t} u_1 (z^t_1(x), z^t_2(x), r, t)
\end{align}
Substitution from the equation \ref{eq:time_der_coupled} to the equation \ref{eq:def_average_v}, we then have the form of the ground truth (target) average velocity field as: 
\begin{equation}
     u_1^{\text{tar}}(z^t_1(x), z^t_2(x), r, t) = u^{\text{ins}}_1(z^t_1(x), z^t_2(x), t ) - (t -r) \frac{d}{dt}  u_1(z^t_1(x), z^t_2(x), r, t) \label{eq:target_def_average_coupled_velocity}
\end{equation}
\subsection{Average Velocity with Multiple PDEs Processes}
\label{sec:avergae_velocity_multiple}
In this section, we provide the analysis of the ground truth average velocity field in the setting of multiple processes $v_i(x,t)$ with $i \in [1, n]$. Derive the same step as Section~\ref{sec:average_velocity_couple}, we have the ground truth (target) average velocity field for process $i-$th as: 
\begin{equation}
     u_i^{\text{tar}}(z^t_1(x),.., z^t_n(x), r, t) = u^{\text{ins}}_i(z^t_1(x),.., z^t_n(x), t ) - (t -r) \frac{d}{dt}  u_i(z^t_1(x),.., z^t_2(x), r, t) \label{eq:target_def_average_multiple_velocity}
\end{equation}
with the time derivative of the average velocity $u_i(z^t_1(x),.., z^t_n(x), r, t)$ as:
\begin{align}
    \frac{d}{dt}  u_i(z^t_1(x),..., z^t_n(x), r, t) &=\partial_t u_i(z^t_1(x), ...,z^t_n(x), r, t) \label{eq:time_der_multiple} \\ \nonumber
    &+ \sum_{j=1}^{n} u^{\text{ins}}_j( z^t_1(x), ...,z^t_n(x), t ) \partial_{z_j^t} u_i (z^t_1(x),..., z^t_n(x), r, t) 
\end{align}

\subsection{Model Instant Velocity Field using Gaussian Process}
In Section~\ref{sec:average_velocity_couple} and Section~\ref{sec:avergae_velocity_multiple}, we provided the analysis to compute the ground truth average velocity field in coupled processes and multiple processes settings. However, we still need to estimate the instant velocity field $u^{\text{ins}}_i$ to compute the target average velocity effectively. In this section, we introduce the framework to compute this instant velocity based on the Gaussian Process~\cite{Rasmussen2006Gaussian} in both coupled processes (Section~\ref{sec:Gaussian_coupled}) and multiple processes settings (Section~\ref{sec:Gaussian_multiple}).

\subsubsection{Coupled PDE setting}
\label{sec:Gaussian_coupled}
In this section, we will provide the analysis in the coupled PDEs setting, where there are only two processes $v_1$ and $v_2$ that interact with each other. Formally, we have the dataset that only includes the initial condition $v_i(x,0)$ and the solution $v_i(x, \tau)$ at time $\tau$. We model the dependence between two random variables $v_1$ and $v_2$ and time $t$ using Gaussian Process, such that, changing the timescale to flow-based $[0, 1]$, we have the $N$ training pairs $\{(0, z^0_1(x)); z^0_2(x)\}$ and $\{(1, z^1_1(x)); z^1_2(x)\}$. The Gaussian process \citep{Rasmussen2006Gaussian} models the joint distribution between any finite number of random variables as a joint Gaussian distribution. We denote $\mathrm{y}$ as $z_2$, $\mathrm{x}$ as vector of  $(t, z_1) \in \mathcal{X} \subset \mathbb{R}^{d+1}$, $\mathrm{y} = f(\mathrm{x}) \in \mathbb{R}^d$, and learn a vector-valued GP~\citep{alvarez2012kernelsvectorvaluedfunctionsreview}, such that:
\begin{equation}
   \mathbf{f} \sim \mathcal{GP}(\mathbf{m}, \mathbf{K})
\end{equation}
where $\mathbf{m} \in \mathbb{R}^d$ is the mean function and $\mathbf{K}$ is a
positive semi-definite Gram 
% matrix-valued 
% function 
such that the value $(\mathbf{K}(\mathrm{x},\mathrm{x}'))_{p,q}$ in the matrix $\mathbf{K}(\mathrm{x}, \mathrm{x}')$ corresponds to
the covariance between the outputs $f_p(\mathrm{x})$ and $f_q(\mathrm{x}')$, which are value of $f(\mathrm{x})$ and $f(\mathrm{x}')$ at dimension $p$ and $q$. For a set of inputs $\mathbf{X}$, the prior distribution over the vector $\mathbf{f}(X)$ is given by 
\begin{align}
    \mathbf{f}(X) \sim \mathcal{N}\big(\mathbf{m}(X), \mathbf{K}(X,X) \big)
\end{align}
The predictive distribution for new vector $\mathrm{x}_{*}$ is 
\begin{align}
    \mathbf{f}(\mathrm{x}_{*}) \mid \mathbf{f}, \mathrm{x}_{*}, X \sim \mathcal{N} \big(\mathbf{f}_{*}(\mathrm{x}_{*}) , \mathbf{K}_{*}(\mathrm{x}_{*}, \mathrm{x}_{*})\big)
\end{align}
where $\mathbf{K}_{\mathrm{x}_*} \in \mathbb{R}^{d \times Nd}$ has entries $(\mathbf{K}(\mathrm{x}_{*}, \mathrm{x}_{j}))_{p,q}$ for $j \in [1,N]$ and $p,q \in [1, d]$ (see Appendix Section~\ref{appendix:GP_details} for implementation details). Thus, given new input $\mathrm{x}_{*} = (t, z_1^*)$, we can predict 
% $z_2(\mathrm{x}_{*})$
$z_2^*$ as
\begin{equation}
    z_2^* = \mathbf{f}_{*} (\mathrm{x}_{*}) \label{eq:infer_z}
\end{equation}
We can also compute the instant velocity of the random variable $z_2$ at time $t$ as 
\begin{equation}
    u^{\text{ins}}_2 (z_2, \mathrm{x}_{*}) = \frac{d}{dt} \Big(  z_2(\mathrm{x}_{*}) \Big) = \frac{d}{dt} \Big(  \mathbf{f}_{*}(\mathrm{x}_{*}) \Big) \label{eq:infer_velocity}
\end{equation}
Note that the derivative of a Gaussian process is another Gaussian process~\cite{Rasmussen2006Gaussian}. Thus we can use GPs to make predictions about the instant velocity based on derivative information. Given predefined covariance function $k(\cdot , \cdot)$, we have covariance between function values and
partial derivatives as
\begin{align}
    \text{cov}\Big(\mathbf{f}_*(\mathrm{x_i}), \frac{\partial \mathbf{f}_*(\mathrm{x_j})}{\partial t}\Big) &= \frac{\partial k(\mathrm{x}_i, \mathrm{x}_j)}{\partial t}  \\ \nonumber
    \text{cov}\Big(\frac{ \partial \mathbf{f}_*(\mathrm{x_i})}{ \partial t}, \frac{\partial \mathbf{f}_*(\mathrm{x_j})}{\partial t}\Big) &= \frac{\partial^2 k(\mathrm{x}_i, \mathrm{x}_j)}{\partial t}
\end{align}

So that, given $\mathrm{x}_* =(t, z_1)$ we can sample the value of $z_2(\mathrm{x}_*)$ and also compute the instant velocity $u^{\text{ins}}_2 (z_2, \mathrm{x}_*)$, a lasting bottleneck is how we can sample $(t, z_1)$. Inspired from Dice Strategy~\citep{xiao2023coupled} to randomly decide the interaction order between each process in a physical system in the current layer. We first sample $t \sim U(0,1)$ from a uniform distribution, and we then deploy the dice strategy to randomly decide the interaction order between processes $v_1$ and $v_2$ with a predefined probability $p \in [0,1]$. If $v_1$ interact to $v_2$, we model the value of $z^t_1$ as the linear interpolant of time with a fixed ending point
\begin{equation}
    z^t_1 = (1-t)z^0_1 + t z^1_1
\end{equation}
Then, we can plug $\mathrm{x}_* =(t, z^t_1)$ into Eq.~\ref{eq:infer_z} and Eq.~\ref{eq:infer_velocity} to compute $z^t_2$ as well as its instant velocity field.
\subsubsection{Multiple PDE setting}
\label{sec:Gaussian_multiple}
In this section, we will provide the analysis of multiple PDE settings, where there are $n$ processes $v_1, ..., v_n$ that interact with each other. To extend the coupled PDE setting to multiple PDE setting, the simple idea is to construct multiple GP mappings: $(t, z_i) \to z_j$ with $i,j \in [1,n]$ and $i \neq j$, and using the Dice Strategy to randomly decide the interaction order between processes. Given each pair $(r, t)$, we randomly decide a driver process $i^*$, which interacts with other processes. We can sample the driven processes as 
\begin{equation}
    z^t_{i^*} = (1-t)z^0_{i^*} + t z^1_{i^*}
\end{equation}
We then use $n-1$ GP processes to model the interaction among individual processes.


\subsection{Training Algorithm with GP modeling}
In this section, we provide our training algorithm using the flow-based model to solve coupled PDE and multiple PDE with the average velocity derived in Section~\ref{sec:average_velocity_couple} and Section~\ref{sec:avergae_velocity_multiple} and $\mathcal{GP}$ modeling. We first denote $Z(t) = (z^t_1, z^t_2, ..., z^t_n)$ and introduce the global average velocity loss for each process $i$ as 
\begin{equation}
    \mathcal{L}^{i}_{\text{global}}(\theta) = \|u^{\theta}_i\big(Z(1), 0,1 \big) - (z^1_i - z^0_i)\|^2_2 \label{eq:global_gp_loss_i}
\end{equation}
Then, for $n$ processes, we have the global loss as
\begin{equation}
    \mathcal{L}_{\text{global}}(\theta) = \sum_{i=1}^{n}  \mathcal{L}^{i}_{\text{global}}(\theta)\label{eq:global_gp_loss}
\end{equation}
We note that the average velocity fields not only need to satisfy the global constraints, but also need to satisfy the local constraints, which are derived from their definitions at each time $t$. Let $u_i^{\text{tar}}(Z(t),r,t)$ be the target average velocity field from Eq.~\ref{eq:target_def_average_coupled_velocity} for coupled PDE setting and Eq.~\ref{eq:target_def_average_multiple_velocity} for multiple PDE setting. We have the local velocity loss for each process $i$ as 
\begin{equation}
    \mathcal{L}^{i}_{\text{local}}(\theta, r,t) = \|u^{\theta}_i\big(Z(t), r,t\big) - \text{sg}\big(u_i^{\text{tar}}(Z(t),r,t)\big) \|^2_2 \label{eq:local_gp_loss_i}
\end{equation}
In the local loss function, a stop-gradient (sg) operation is applied on the target average velocity fields $u_i^{\text{tar}}(Z(t),r,t)$ following the common practice~\citep{geng2025mean} to avoid double backpropagation. For $n$ processes, we have the combined local loss as
\begin{equation}
    \mathcal{L}_{\text{local}}(\theta, r,t) = \sum_{i=1}^{n}\mathcal{L}^{i}_{\text{local}}(\theta, r,t) \label{eq:local_gp_loss}
\end{equation}
The detailed training and sampling algorithm are presented in Alg.~\ref{alg:training} and Alg.~\ref{alg:sampling}. 
\begin{algorithm}
\caption{Training Algorithm with GP modeling}
\label{alg:training}
\begin{algorithmic}[1]
\STATE \textbf{Input:} Dataset $\mathcal{D}$, number of epochs $n_0$.
\STATE \textbf{Output:} The flow model weight $\theta$.
\FOR{$i = 1$ to $n_0$}
    \STATE  \texttt{t,r = sample\_t\_r()}
    \STATE Sample $Z(t) = (z^t_1,...,z^t_n)$ using $\mathcal{GP}$ (Section~\ref{sec:avergae_velocity_multiple}) 
    \STATE Compute local loss $\mathcal{L}_{\text{local}}(\theta,r,t)$ as Eq.~\ref{eq:local_gp_loss}
    \STATE Update $\theta$ via minimizing the local loss $\mathcal{L}_{\text{local}}(\theta,r,t)$
    \STATE Compute global loss $\mathcal{L}_{\text{global}}(\theta)$ as Eq.~\ref{eq:global_gp_loss}
    \STATE Update $\theta$ via minimizing the global loss $\mathcal{L}_{\text{global}}(\theta)$
\ENDFOR
\STATE \textbf{Return} the weight $\theta$
\end{algorithmic}
\end{algorithm}

\begin{algorithm}
\caption{One-Step Sampling Algorithm}
\label{alg:sampling}
\begin{algorithmic}[1]
\STATE \textbf{Input:} Test point $Z(1) = (z^1_1, ..., z^1_n)$, pretrained weight $\theta$
\FOR{$i = 1 $ to $n$}
   \STATE $z^0_i = z^1_i - u^{\theta}_{i}(Z(1),0,1)$
\ENDFOR
\STATE $Z(0) = (z^0_1, ..., z^0_n)$
\STATE Return $Z(0)$
\end{algorithmic}
\end{algorithm}


\section{Theoretical Analysis}
In this section, we present our main theoretical results. We begin with problem settings and the key notation. We then establish a generalization error bound for our method that characterizes how the local loss $\mathcal{L}_{\text{local}}$ and global loss $\mathcal{L}_{\text{global}}$ affect the model's performance. This provides a principled foundation for understanding coupled flow models for solving coupled PDEs.
\subsection{Problem formulation and notations}
\label{sec:notation}
For consistency with the literature, let the source distribution with each process $i \in[1,n]$ be $\pi^{(i)}_{0}$, the joint source distribution among processes be $\pi_{0}$. Denote the target distribution with each process $i$ as $\pi^{(i)}_{\tau}$ and the joint target distribution among processes be $\pi_{\tau}$, the predictive distribution as $\hat{\pi}^{(i)}_{\tau}$ and $\hat{\pi}_{\tau}$, respectively. Let $W_2$ be the Wasserstein-2 distance, we have 
\begin{equation}
    \sum_{i=1}^{n}W^2_2(\pi^{(i)}_{\tau}, \hat{\pi}^{(i)}_{\tau}) \leq W^2_2 (\pi_\tau, \hat{\pi}_{\tau})
\end{equation}
As discussed in Section~\ref{sec:method}, we have the approximate average velocity field for each process as $u^{\theta}_{i}(Z_t, r, t)$. We define the approximate instant velocity field of each processes $i$ by
\begin{equation}
    \mu^{\theta}_{i}(t, Z_t) = \lim_{\Delta t \to 0}  u^{\theta}_{i}(Z_t, t, t+ \Delta t) = u^{\theta}_{i}(Z_t, t,t)
\end{equation}
We then define the approximate flow map among processes as $M^{\theta}(t, Z_t) = [\mu^{\theta}_1(t, Z_t), ..., \mu^{\theta}_{n}(t, Z_t)]$. The results of this section will be proved under the following common assumptions~\citep{benton2024error}. 

\begin{assumption}[Existence and uniqueness of smooth flows]
\label{assum:existence}
For each $\mathrm{z} \in \mathbb{R}^{n \times d}$ and $s\in [0, \tau]$, there exist unique flows $(Y^{\mathrm{z}}_{s,t})_{t\in [s, \tau]}$ and $(Z^{\mathrm{z}}_{s,t})_{t\in [s, \tau]}$ starting at $\mathrm{z}$ with flow map $M^{\text{tar}}(t, \mathrm{z})$ and $M^{\theta}(t, \mathrm{z})$ respectively and continuously differentiable in the input $\mathrm{z},s,t$. 
\end{assumption}
\begin{assumption}[Regularity of approximate flow map]
\label{assum:lipschitz}
The approximate flow map $M^{\theta}(t, \mathrm{z})$ is differentiable in both inputs. Also, for each $t \in (0, \tau)$, there is a constant $L_t$ such that $M^{\theta}(t, \mathrm{z})$ is $L_t$-Lipschitz in $\mathrm{z}$
\end{assumption}

\subsection{Main Result}
\label{sec:theoretical_results}
Our first main result is the following, which bounds the error of the flow-matching procedure in the Wasserstein-2 distance in terms of the instantaneous local loss and the Lipschitz constant of the approximate flow map $M^{\theta}(t, \mathrm{z})$.

\begin{theorem}
\label{thm:main_bound}
Suppose we have $\pi^{(i)}_{\tau}$ and $\hat{\pi}^{(i)}_{\tau}$ as defined in section~\ref{sec:notation} that $Z$ is the flow starting in $\pi_0$ with flow map $M^{\theta}(t, \mathrm{z})$. Under Assumptions \ref{assum:lipschitz} and \ref{assum:existence}, we have 
$$ \sum_{i=1}^{n} W_2^2(\pi^{(i)}_{\tau}, \hat{\pi}^{(i)}_{\tau}) \leq \mathcal{L}_{\text{global}}(\theta) \leq K^2_{0,\tau} \mathbb{E}_{t \sim U(0,\tau)}[\mathcal{L}_{\text{local}}(\theta,t,t)]$$ 
where $K_{0,\tau} = \exp \Big(\int_{0}^{\tau} L_t dt \Big)$ reflect the smoothness of approximate flow map. The detailed proof provided in Appendix~\ref{appendix:Proof_mean_bound}.
\end{theorem}
This result reveals how the distributional differences between the ground truths and predictions among individual processes can be upper bounded by the global loss $\mathcal{L}_{\text{global}}$, the local loss $\mathcal{L}_{\text{local}}$ and the smoothness controlling term $K_{0, \tau}$, which is the integral of the Lipschitz constant of the approximate flow map $L_t$. We next, extend the setting to the rollout prediction, where the flow model trained on coupling $(\pi_0, \pi_{\tau})$, when predict on time step $\tau + \Delta \tau$ with $\Delta\tau \geq 0$. Using Theorem~\ref{thm:main_bound}, we have the following result.

\begin{lemma}[Error Bound in Rollout Prediction]
\label{lem:err_rollot}
Suppose we have $\pi^{(i)}_{\tau + \Delta \tau}$ and $\hat{\pi}^{(i)}_{\tau+ \Delta \tau}$ as the ground truth and predictive distribution respectively, that $Z$ is the flow starting in $\pi_0$ with flow map $M^{\theta}(t, \mathrm{z})$. Extending Assumptions \ref{assum:lipschitz} and \ref{assum:existence} to time-step $\tau + \Delta \tau$, we have 
$$ \sum_{i=1}^{n} W_2^2(\pi^{(i)}_{\tau+ \Delta \tau}, \hat{\pi}^{(i)}_{\tau + \Delta \tau}) \ \leq K^2_{0,\tau+ \Delta \tau} \Big(\mathbb{E}_{t \sim U(0,\tau)}[\mathcal{L}_{\text{local}}(\theta,t,t)] +  \mathbb{E}_{t \sim U(\tau,\tau+\Delta \tau)}[\mathcal{L}_{\text{local}}(\theta,t,t)]\Big)$$  
\end{lemma}


\section{Empirical Analysis}
\subsection{Implementation details}
\textbf{Benchmarks}:  Following the prior practice \citep{sun2025compolunifiedneuraloperator}, we evaluated our method on predicting coupled multi-physics PDEs systems using five benchmarks: 1-D Lotka-Volterra equations \citep{Murray2002} (2 processes), 1-D Belousov-Zhabotinsky equations \citep{Petrov1993} (3 processes), 2-D Gray-
Scott equations \citep{Pearson_1993} (2 processes), 2-D multiphase flow of oil and water in geological
storage \citep{bear2010modeling} \citep{hashemi2021pore} \cite{aboukassem2013petroleum} (2 processes), and 2-D Thermal-Hydro-Mechanical problem \citep{Gao2020} (5 processes). The training datasets were stimulated with numerical solvers using 256-point meshes for 1-D cases and $64 \times 64$ meshes for 2-D cases. As prior practice \citep{sun2025compolunifiedneuraloperator}, we also trained the models on two dataset sizes ($512$ and $1024$ samples) and evaluated performance on an independent test set of $200$ samples. Coefficients and boundary
conditions remained constant to isolate the effects of varying initial conditions.

\textbf{Competing baselines}: We evaluated our proposed method again state-of-the-art neural operators baselines: COMPOL \citep{sun2025compolunifiedneuraloperator}, FNO \citep{li2021fourier}, UFNO \citep{wen2022ufnoenhancedfourier}, Transolver \citep{wu2024transolverfasttransformersolver}, DeepOnet \citep{lu2021learning}, and CMWNO \citep{xiao2023coupled}. COMPOL and CMWNO are designed for coupled PDE problems, while other baselines are not inherently designed for multi-physics modeling. Therefore, following the practice~\cite{sun2025compolunifiedneuraloperator}, we adapt them by concatenating process inputs and outputs across their respective channels. We also adapt the prior experiment setting \citep{sun2025compolunifiedneuraloperator}, \cite{xiao2023coupled}, with Adam optimizer (learning rate $=0.001$). Except for CMWNO, which employed a step scheduler (as in its original
implementation), all models used a cosine annealing schedule. DeepONet is trained with $10000$ epochs, while other baselines trained with $500$ epochs. 

\subsection{Analysis of Predictive Performance}

\begin{table}[!htbp]
\centering
\renewcommand{\arraystretch}{1.2}
\setlength{\tabcolsep}{0.35em}
\resizebox{\columnwidth}{!}{%
\begin{tabular}{@{}lc|ccccc@{}}
\hline
\textbf{Method} & \textit{Train Samples} & \textit{Lotka-V.} & \textit{Belousov-Z.} & \textit{Grey-Scott} & \textit{Multiphase} & \textit{THM} \\
\hline
\multirow{2}{*}{$\text{FNO}_{c}$} & 512 & 12.29 & 0.18 & 0.88 & 2.29 & 0.15 \\
& 1024 & 9.55 & 0.16 & 0.55 & 7.98 & 0.11 \\
\hline
\multirow{2}{*}{CMWNO} & 512 & 27.21 & 0.64 & 0.83 & 9.54 & 3.30 \\
& 1024 & 16.55 & 0.51 & 0.63 & 6.86 & 2.72 \\
\hline
\multirow{2}{*}{$\text{UFNO}_{c}$} & 512 & 11.63 & 0.20 & 0.55 & 4.48 & 0.13 \\
& 1024 & 8.73 & 0.14 & 0.36 & 7.61 & 0.08 \\
\hline
\multirow{2}{*}{$\text{Transolver}_{c}$} & 512 & 18.19 & 5.97 & 0.54 & 1.28 & 0.18 \\
& 1024 & 17.67 & 1.88 & 0.49 & 0.88 & 0.16 \\
\hline
\multirow{2}{*}{$\text{DiffusionPDE}_{c}$} & 512 & 472.41 & 308.81 & 31.67 & 12.84 & 5.62 \\
& 1024 & 455.28 & 172.74 &  & 14.33 & 5.62 \\
\hline
\multirow{2}{*}{$\text{DeepONet}_{c}$} & 512 & 55.46 & 1.11 & 31.30 & 31.20 & 31.36 \\
& 1024 & 54.83 & 0.99 & 30.99 & 28.19 & 31.50 \\
\hline
\multirow{2}{*}{COMPOL-RNN} & 512 & 10.43 & 0.17 & 0.66 & 2.21 & 0.14 \\
& 1024 & 8.02 & 0.15 & 0.50 & 6.44 & 0.09 \\
\hline
\multirow{2}{*}{COMPOL-ATN} & 512 & 11.32 & 0.18 & 0.61 & 2.17 & 0.14 \\
& 1024 & 8.58 & 0.15 & 0.48 & 6.76 & 0.10 \\
\hline
\multirow{2}{*}{Our-GP} & 512 & \textbf{2.16} & \textbf{0.15} & \textbf{0.19} & \textbf{0.52} & \textbf{0.07} \\
& 1024 & \textbf{2.01} & \textbf{0.09} & \textbf{0.15} & \textbf{0.13} & \textbf{0.06} \\
\hline
\end{tabular}}
\vspace{1em}
\caption{\small Comparison of relative $L_2$ prediction errors ($\times 10^{-2}$). Bold values indicate the best performance for each case. }
\label{tab:main_results}
\end{table}



\section{Discussion}
[TODO]
\section{Conclusion}
[TODO]


%%%%%%%%%%%%%%%%%%%%%%%%%%%%%%%%%%%%%%%%%%%%%%%%%%%%%%%%%%%%
\newpage
\appendix

\section{GP Modeling implementation details}
\label{appendix:GP_details}
In this section, we provide the details about $\mathcal{GP}$ modeling for the interpolant of the flow-matching path. As discussed in Section~\ref{sec:Gaussian_coupled}, we model the dependence between two random variable $z_1$ and $z_2$ and time $t$ using Gaussian Processes. The idea here is to model $z^t_1$ as the function of time $t$ and $z^t_2$, and otherwise. That way, we can capture the correlation information among individual processes by randomly deciding the interaction order at each time-step $t$ following the prior practice \citep{xiao2023coupled}. For example, at time-step $t$, we decided the variable $z^t_2$ will impact on the variable $z^t_1$. Let $\mathrm{x} = (t, z^t_2) \in \mathcal{X} \subset \mathbb{R}^{d+1}$ and $\mathrm{y} = z^t_1 = \mathrm{f}(\mathrm{x}) \in \mathbb{R}^{d}$, we model the function $\mathrm{f}$ as a vector-valued GP \citep{alvarez2012kernelsvectorvaluedfunctionsreview}, such that 
\begin{equation}
   \mathbf{f} \sim \mathcal{GP}(\mathbf{m}, \mathbf{K})
\end{equation}

Given the input set as the pair data at beginning time $t=0$ and ending time $t=1$ that $X = \{\mathrm{x}_0; \mathrm{x}_1  \}=\{ (0, z^0_2);  (1, z^1_2)\}$, the prior distribution over the function $\mathrm{f}(X)$ is given by 
\begin{equation}
    \mathrm{f}(X) \sim \mathcal{N}(\mathbf{m}(X), \mathbf{K}(X,X))
\end{equation}
Given time step $t$, we directly sample $z^t_2$ as the linear interpolant of time,
\begin{equation}
    z^t_2 = (1-t) z^0_2 + t z^1_2
\end{equation}
Thus, we can construct the new vector $\mathrm{x}_* = (t, z^t_2)$, with the noise-free $\mathcal{GP}$ modeling, we have the predictive distribution as  
\begin{equation}
    \mathrm{f}(\mathrm{x}_*) \mid \mathrm{f}, \mathrm{x}_*,X \sim \mathcal{N}(\mathbf{f}_*(x_*), \mathbf{K}_*(\mathrm{x}_*, \mathrm{x}_*))
\end{equation}
where 
\begin{align}
    \mathbf{f}_{*}(\mathrm{x}_*) &= \mathbf{K}^{T}_{\mathrm{x}_{*}} \big(\mathbf{K}(X,X) )^{-1} \mathbf{y} \\ 
    \mathbf{K}_{*}(\mathrm{x}_{*}, \mathrm{x}_{*}) &= \mathbf{K}(\mathrm{x}_{*}, \mathrm{x}_{*}) - \mathbf{K}_{\mathrm{x}_{*}}\big(\mathbf{K}(X,X)\big)^{-1}\mathbf{K}^{T}_{\mathrm{x}_{*}}
\end{align}
with provided data point $\mathrm{y}_0 = z^0_1$ and $\mathrm{y}_1 = z^1_1$, we have the label vector $\mathbf{y} = [\mathrm{y}_0  \: \mathrm{y}_1]^T \in \mathbb{R}^{2d}$. To encode the correlation information among the output dimensions, we use separable kernels \citep{alvarez2012kernelsvectorvaluedfunctionsreview} as the output kernel function, such that 
\begin{equation}
    \mathbf{K}(\mathrm{x}_i, \mathrm{x}_j)  = k(\mathrm{x}_i, \mathrm{x}_j) B
\end{equation} 
where $k (\mathrm{x}_i, \mathrm{x}_j)$ are scalar kernels on $\mathcal{X} \times \mathcal{X}$ and $B\in \mathbb{R}^{d \times d}$ is the symmetric and positive semi-definite matrix. Thus, we have the matrix expression as 
\begin{equation}
    \mathbf{K}(X, X) = \mathbf{k}_{\mathrm{x}}(X,X)  \otimes  \mathbf{B} 
\end{equation}
where 
\begin{equation}
    \mathbf{k}_{\mathrm{x}}(X,X) = \begin{bmatrix}
k(\mathrm{x}_0, \mathrm{x}_0) & k(\mathrm{x}_0, \mathrm{x}_1) \\
k(\mathrm{x}_1, \mathrm{x}_0) & k(\mathrm{x}_1, \mathrm{x}_1)
\end{bmatrix}
\end{equation}
Similarly, the query-to-train covariance $\mathbf{K}^{T}_{\mathrm{x}_{*}}$ is
\begin{equation}
    \mathbf{K}^{T}_{\mathrm{x}_{*}} = [  k(\mathrm{x_*}, \mathrm{x}_0)B  \quad k(\mathrm{x}_{*}, \mathrm{x}_{1})B] = \mathbf{k}^{T}_{\mathrm{x}_{*}} \otimes B
\end{equation}
where 
\begin{equation}
   \mathbf{k}^{T}_{\mathrm{x}_{*}} =  [  k(\mathrm{x_*}, \mathrm{x}_0)  \quad k(\mathrm{x}_{*}, \mathrm{x}_{1})]
\end{equation}
Thus, we have the mean function as $\mathbf{f}_{*}(\mathrm{x}_*)$ as 
\begin{align}
    \mathbf{f}_{*}(\mathrm{x}_*) &= \mathbf{K}^{T}_{\mathrm{x}_{*}} \big(\mathbf{K}(X,X) )^{-1} \mathbf{y} \\
    &= \big(\mathbf{k}^{T}_{\mathrm{x}_{*}} \otimes \mathbf{B} \big)\big(\mathbf{k}_{\mathrm{x}}(X,X)^{-1}  \otimes  \mathbf{B}^{-1} \big) \mathbf{y} \\
    &= \Big(\mathbf{k}^{T}_{\mathrm{x}_{*}} \mathbf{k}_{\mathrm{x}}(X,X)^{-1}\Big) \otimes \Big(\mathbf{B}  \mathbf{B}^{-1}\Big) \mathbf{y} \\
    &= \Big( \big(\mathbf{k}^{T}_{\mathrm{x}_{*}} \mathbf{k}_{\mathrm{x}}(X,X)^{-1}\big) \otimes I_{d} \Big)\mathbf{y} \label{eq:mean_function_closed_form}
\end{align}
Using Eq.~\ref{eq:mean_function_closed_form}, we have the closed form of the mean function when predicting for the new input $\mathrm{x}_{*}$ without the requirement for the pre-trained matrix $B$. For practical implementation, we define the input covariance function as the squared-exponential covariance function 
\begin{equation}
    k(\mathrm{x}_i, \mathrm{x}_j) = \sigma^2_{f} \exp\Big( - \frac{1}{2 \ell^2} \| \mathrm{x}_i - \mathrm{x}_j\|^2_2\Big)
\end{equation}
where $\ell$ is the length-scale and $\sigma^2_f$ is the signal variance. 


\section{Proof of the Theorem~\ref{thm:main_bound}}
\label{appendix:Proof_mean_bound}
In this section, we provide the detailed proof of the theorem~\ref{thm:main_bound}. As discussed in Section~\ref{sec:notation}, for flow gap $M(t, Z_t)$, we get the standard ODE as 
\begin{equation}
    M(t, Z ) = \frac{d}{dt}\Big(Z(t) \Big)
\end{equation}
As Section~\ref{sec:notation}, we have the ground truth distribution at the solution space be $\pi_\tau$, and the predicted distribution be $\hat{\pi}_\tau$. Then, 
\begin{equation}
    \sum_{i=1}^{n} W^2_2(\hat{\pi}^{(i)}_\tau, \pi^{(i)}_\tau) \leq W^2_2(\hat{\pi}_\tau, \pi_\tau)  \leq \mathbb{E} \Big[\| Z_\tau - Y_\tau \|^2_2 \Big] = \mathcal{L}_{\text{global}}(\theta)
\end{equation}
where $Y(t)$ is the ground truth flow variable with ground truth velocity field as $\dot{Y}(t)$ and  $Z: \mathbb{R}^{n \times d} \times [0,T]^2 \to \mathbb{R}^{n \times d}$ be a continuous solution of the integral equation 
\begin{equation}
    Z^{\mathrm{z}}_{s,t} = \mathrm{z} + \int_{s}^{t} M^{\theta}(r, Z^{\mathrm{z}}_{s,r}) dr \label{eq:define_learned_flow}
\end{equation}
Using the Alekseev-Gröbner~\citep{grobner1960lie} formula, we have the following. 
\begin{equation}
    Z^{Y_0}_{0,T} - Y_T = \int_{0}^{T} \Big(\nabla_{Y_r} Z^{Y_r}_{r,T}\Big) \Big(M^{\theta}(r,Y_r) - \dot{Y}_r \Big) dr  \label{eq:Alek_Gro_formula}
\end{equation}
Using Eq.~\ref{eq:define_learned_flow}, we can write 
\begin{equation}
    \nabla_{\mathrm{z}} Z^{\mathrm{z}}_{s,t} = I + \int_{s}^{t} \nabla_{\mathrm{z}} M^{\theta}(r, Z^{\mathrm{z}}_{s,r}) \nabla_{\mathrm{z}} Z^{\mathrm{z}}_{s,r} dr
\end{equation}
It shows that 
\begin{equation}
    \frac{\partial}{\partial t} \|\nabla_{\mathrm{z}}  Z^{\mathrm{z}}_{s,t}\|_{\text{op}} \leq \| \frac{\partial}{\partial t} \nabla_{\mathrm{z}}  Z^{\mathrm{z}}_{s,t}\|_{\text{op}} = \| \nabla_{\mathrm{z}} M^{\theta}(r, Z^{\mathrm{z}}_{s,t})  \nabla_{\mathrm{z}} Z^{\mathrm{z}}_{s,t}\|_{\text{op}}
\end{equation}
where $\|\cdot \|_{\text{op}}$ denotes the operator norm with respect to the $\ell_2$ metric. With the assumption that, $\| \nabla_{\mathrm{z}} M^{\theta}(r, Z^{\mathrm{z}}_{s,t})\| \leq L_t$, we have 
\begin{equation}
   \frac{\partial}{\partial t} \|\nabla_{\mathrm{z}}  Z^{\mathrm{z}}_{s,t}\|_{\text{op}} \leq L_t \|\nabla_{\mathrm{z}}  Z^{\mathrm{z}}_{s,t}\|_{\text{op}} \label{eq:norm_leq_Lt}
\end{equation}
Using Grönwall’s lemma, we have 
\begin{equation}
    \|\nabla_{\mathrm{z}}  Z^{\mathrm{z}}_{s,t}\|_{\text{op}} \leq \exp\Big(\int^t_s L_r dr \Big) \leq \exp\Big(\int^\tau_0 L_r dr \Big)
\end{equation}
Denote $K_{0, \tau}:= \exp\Big(\int^\tau_0 L_r dr \Big)$ then plug into Eq.~\ref{eq:Alek_Gro_formula} at $T=\tau$, we get
\begin{equation}
    \|Z^{\mathrm{z}}_{0,\tau}  - Y^{\mathrm{z}}_{0,\tau}\|_2 \leq K_{0, \tau} \int_{0}^{\tau} \|M^{\theta}(r, Y_r) - \dot{Y}_{r} \|_2 dr
\end{equation}
Taking expectation over $\mathrm{z} \sim \pi_0$, we have 
\begin{align}
    \mathbb{E}\big[ \| Z^{\mathrm{z}}_{0, \tau}  - Y^{\mathrm{z}}_{0, \tau}\|_2^2\big] &\leq K^2_{0, \tau} \cdot\mathbb{E} \Big[\Big( \int_{0}^{\tau} \|M^{\theta}(r, Y_r) - \dot{Y}_{r} \|_2 dr\Big)^2 \Big]  \\
    &\leq K_{0, \tau}^2 \int_{0}^\tau \mathbb{E} \Big[\|M^{\theta}(r, Y_r) - \dot{Y}_{r} \|^2_2 \Big] dr \label{eq:error_bound_raw}
\end{align}
As discussed in Section~\ref{sec:notation}, we note that the instantaneous velocity field, $\mu^{\theta}_i(r, Y_r)$, represents the limiting value of the time-averaged velocity field as the duration of the observation interval, $\Delta t$, vanishes, such that 
\begin{equation}
    \mu^{\theta}_i(r, Y_r) =\lim_{\Delta t \to 0} u^{\theta}_i (Y_r, r, r + \Delta t) = u^{\theta}_i(Y_r, r, r )  
\end{equation}
Thus, we can write 
\begin{align}
    \mathbb{E} \Big[\|M^{\theta}(r, Y_r) - \dot{Y}_{r} \|^2_2 \Big] &= \sum_{i=1}^{n} \mathbb{E} \Big[\|\mu_i^{\theta}(r, Y_r) -  \mu_i^{\text{tar}}(r, Y_r)\|^2_2 \Big] \\
    &= \mathcal{L}_{\text{local}}(\theta, r, r) \label{eq:relate_to_local_loss}
\end{align}

Finally, using Eq.~\ref{eq:relate_to_local_loss} and  Eq.~\ref{eq:error_bound_raw}, we have the theorem as
\begin{equation}
    \sum_{i=1}^{n} W_2^2(\pi^{(i)}_{\tau}, \hat{\pi}^{(i)}_{\tau}) \leq \mathcal{L}_{\text{global}}(\theta) \leq K^2_{0,\tau} \mathbb{E}_{t \sim U(0,\tau)}[\mathcal{L}_{\text{local}}(\theta,t,t)]
\end{equation}

\section{Smoothness of velocity fields across design choices}
\label{appedix:Smoothness_design_choices}
As discussed in Section~\ref{sec:theoretical_results}, the bound in the Theorem is exponentially dependent on the Lipchitz constant of the learned $M^{\theta}(t, Z)$. In this section, we control the corresponding Lipschitz constant of the chosen target velocity field $\dot{Z}_t$ across different choices of modeling. Follow practice~\citep{benton2024error}, we refer this to controlling the Lipschitz constant of $M^{\theta}(t, Z)$ directly since this is determined by our choice of $\mathcal{M}$.
\subsection{Smoothness of velocity fields with linear interpolant of time}
\label{appendix:smoothness_linear}
We first provide the analysis for the naive setting, where the intermediate latent feature samples form the linear interpolant of time. We have 
\begin{align}
    Z_t&= (1-t) Z_0 + t Z_1 \\
    \dot{Z}_t &= Z_1- Z_0 
\end{align}
Thus, we have the target flow map as 
\begin{align}
    M^{\text{tar}}(t, \mathrm{z}) &= \mathbb{E}[\dot{Z}_t \mid Z_t = \mathrm{z}] = \frac{1}{t} \Big(\mathrm{z} - \mathbb{E}[Z_0 \mid Z_t = \mathrm{z}] \Big)  \\
     \nabla_{\mathrm{z}}M^{\text{tar}}(t, \mathrm{z}) &= \frac{1}{t} \Big(I - \nabla_{\mathrm{z}} \mathbb{E}[Z_0 \mid Z_t = \mathrm{z}] \Big)
\end{align}
We can then compute $K^{\text{linear}}$ as 
\begin{equation}
    K^{\text{linear}} = \exp \Big(\int_{0}^{1} L^*_t dt \Big) =  \exp \Big(\int_{0}^{1}  \sup_{\mathrm{z}} \Big\|\frac{1}{t}\Big(I - \nabla_{\mathrm{z}} \mathbb{E}[Z_0 \mid Z_t = \mathrm{z}] \Big)  \Big\| dt \Big)
\end{equation}

\subsection{Smoothness of velocity fields with GP modeling}
\label{appendix:smoothness_gp}
In this section, we provide the GP modeling, where interaction information among processes is encoded via Gaussian processes and the Dice Strategy~\cite{xiao2023coupled}. We can write the intermediate latent feature at time-step $t$ as 
\begin{align}
    Z_t &= \omega_{1}(t, Z_0, Z_1) Z_0 + \omega_2(t, Z_0, Z_1) Z_1 \\
    \dot{Z}_{t} &=  \dot{\omega}_{1}(t, Z_0, Z_1) Z_0 + \dot{\omega}_2(t, Z_0, Z_1) Z_1
\end{align}
Thus, we have the target flow map as 
\begin{align}
    M^{\text{tar}}(t, \mathrm{z}) &= \mathbb{E}[\dot{Z}_t \mid Z_t = \mathrm{z}] \\
    \nabla_{\mathrm{z}}M^{\text{tar}}(t, \mathrm{z}) &= \nabla_{\mathrm{z}} \mathbb{E}[\dot{Z}_t \mid Z_t = \mathrm{z}] = \mathbb{E}[\nabla_{\mathrm{z}}  \dot{Z}_t \mid Z_t = \mathrm{z}]
\end{align}
We have 
\begin{equation}
    L^*_t = \sup_{\mathrm{z}} \|\mathbb{E}[\nabla_{\mathrm{z}}  \dot{Z}_t \mid Z_t = \mathrm{z}] \|
\end{equation}
We then can compute $K^{\text{gp}}$ as 
\begin{equation}
    K^{\text{gp}} = \exp \Big(\int_{0}^{1} L^*_t dt \Big) = \exp \Big(\int_{0}^{1}  \sup_{\mathrm{z}} \|\mathbb{E}[\nabla_{\mathrm{z}}  \dot{Z}_t \mid Z_t = \mathrm{z}] \| dt \Big)
\end{equation}
\section{Additional Experiment Results and Ablation Studies}
\label{appendix:Experiment_detailed_results}
In this section, we provide the detailed and additional experiment results and ablation studies. 



\begin{table}[!htbp]
\centering
\renewcommand{\arraystretch}{1.2}
\setlength{\tabcolsep}{0.35em}
\resizebox{\columnwidth}{!}{%
\begin{tabular}{@{}lcccccc@{}}
\hline
\textbf{Method} & \textit{Train Samples} & \textit{Lotka-V.} & \textit{Belousov-Z.} & \textit{Grey-Scott} & \textit{Multiphase} & \textit{THM} \\
\hline
\multirow{2}{*}{FNO} & 512 & 0.1229 $\pm$4e-5 & 0.0018 $\pm$3e-5 & 0.0088 $\pm$4e-5 & 0.0229 $\pm$2e-4 & 0.0015 $\pm$2e-5 \\
& 1024 & 0.0955 $\pm$6e-5 & 0.0016 $\pm$3e-5 & 0.0055 $\pm$3e-6 & 0.0798 $\pm$1e-5 & 0.0011 $\pm$1e-6 \\
\hline
\multirow{2}{*}{CMWNO} & 512 & 0.2721 $\pm$3e-5 & 0.0064 $\pm$2e-5 & 0.0083 $\pm$5e-4 & 0.0954 $\pm$2e-5 & 0.0330 $\pm$4e-5 \\
& 1024 & 0.1655 $\pm$1e-5 & 0.0051 $\pm$5e-5 & 0.0063 $\pm$3e-5 & 0.0686 $\pm$7e-4 & 0.0272 $\pm$5e-5 \\
\hline
\multirow{2}{*}{UFNO} & 512 & 0.1163 $\pm$1e-4 & 0.0020 $\pm$6e-6 & 0.0055 $\pm$5e-5 & 0.0448 $\pm$1e-5 & 0.0013 $\pm$3e-5 \\
& 1024 & 0.0873 $\pm$8e-5 & 0.0014 $\pm$2e-5 & 0.0036 $\pm$6e-6 & 0.0761 $\pm$1e-5 & 0.0008 $\pm$4e-5 \\
\hline
\multirow{2}{*}{Transolver} & 512 & 0.1819 $\pm$1e-5 & 0.0597 $\pm$7e-6 & 0.0054 $\pm$9e-6 & 0.0128 $\pm$4e-5 & 0.0018 $\pm$1e-5 \\
& 1024 & 0.1767 $\pm$1e-5 & 0.0188 $\pm$2e-5 & 0.0049 $\pm$2e-5 & 0.0088 $\pm$3e-5 & 0.0016 $\pm$5e-5 \\
\hline
\multirow{2}{*}{DeepONet} & 512 & 0.5546 $\pm$7e-5 & 0.0111 $\pm$1e-5 & 0.3130 $\pm$1e-5 & 0.3120 $\pm$1e-5 & 0.3136 $\pm$1e-5 \\
& 1024 & 0.5483 $\pm$6e-5 & 0.0099 $\pm$1e-5 & 0.3099 $\pm$1e-5 & 0.2819 $\pm$7e-5 & 0.3150 $\pm$1e-5 \\
\hline
\multirow{2}{*}{COMPOL-RNN} & 512 & 0.1043 $\pm$2e-5 & 0.0017 $\pm$2e-5 & 0.0066 $\pm$1e-4 & 0.0221 $\pm$1e-5 & 0.0014 $\pm$3e-5 \\
& 1024 & 0.0802 $\pm$4e-5 & 0.0015 $\pm$8e-6 & 0.0050 $\pm$4e-5 & 0.0644 $\pm$3e-5 & 0.0009 $\pm$4e-5 \\
\hline
\multirow{2}{*}{COMPOL-ATN} & 512 & 0.1132 $\pm$1e-4 & 0.0018 $\pm$4e-5 & 0.0061 $\pm$3e-5 & 0.0217 $\pm$3e-5 & 0.0014 $\pm$3e-5 \\
& 1024 & 0.0858 $\pm$5e-5 & 0.0015 $\pm$3e-5 & 0.0048 $\pm$2e-5 & 0.0676 $\pm$3e-5 & 0.0010 $\pm$2e-5 \\
\hline
\multirow{2}{*}{Our-GP} & 512 & \textbf{0.0216 $\pm$2e-5} & \textbf{0.0015 $\pm$4e-5} & \textbf{0.0019 $\pm$3e-5} & \textbf{0.0052 $\pm$3e-5} & \textbf{0.0007 $\pm$2e-5} \\
& 1024 & \textbf{0.0201 $\pm$3e-5} & \textbf{0.0009 $\pm$3e-5} & \textbf{0.0015 $\pm$5e-5} & \textbf{0.0013 $\pm$1e-5} & \textbf{0.0006 $\pm$2e-5} \\
\hline
\end{tabular}}
\vspace{1em}
\caption{\small Comparison of relative $L_2$ prediction errors (mean $\pm$ std where std is rounding leftover, or $1\times10^{-5}$ when exact zero). Bold values indicate the best performance for each case.}
\label{tab:nrmse}
\end{table}


\newpage
\input{checklist.tex}

\bibliographystyle{plainnat}
\bibliography{references}
\end{document}